\documentclass[11pt,a4paper]{article}
\usepackage[T1]{fontenc}
\usepackage[utf8]{inputenc}
\usepackage{lmodern}
\usepackage[margin=25mm]{geometry}
\usepackage{microtype}
\usepackage[hidelinks]{hyperref}
\hypersetup{
  pdftitle={Understanding the Mundici functor (my attempt!).},
  pdfauthor={Giuseppe Metere},
  pdfsubject={Online Algebra, Logic and Topology Seminar}
}
\setlength{\parindent}{0pt}

\begin{document}
\begin{center}
{\Large\bfseries Online Algebra, Logic and Topology Seminar}\\[4pt]
{\large (CMUC, University of Coimbra)}
\end{center}

\vspace{1em}
\textbf{Speaker:} Giuseppe Metere (DeFENS, Universit\`a degli Studi di Milano, Italy)\\
\textbf{Date:} October 20th, 2026\\
\textbf{Time:} 15:00 (Europe/Lisbon)\\
\textbf{Place:} Google Meet

\vspace{1.5em}
\begin{center}
{\Large\bfseries Understanding the Mundici functor (my attempt!).}
\end{center}
\vspace{1em}

\textbf{Abstract.} Mundici's equivalence between MV-algebras and lattice-ordered abelian groups with strong unit is usually presented through the construction of the monoid of good sequences. While effective, this construction can appear rather ad hoc from a categorical point of view.

\medskip
In this talk, we reinterpret the first step of Mundici's construction through a universal property. This yields a conceptual characterisation of the monoid of good sequences and places the construction in a more formal setting.

\medskip
We then explain how this universal property leads naturally to a profunctorial perspective on the Mundici functor, and we discuss how the combinatorics of good sequences can be understood through a normalization process reminiscent of polynomial multiplication.

\medskip
This is joint work with Sara Lapenta and Luca Spada.
\end{document}
